\documentclass[10pt,a4paper]{amsart}
\usepackage{amsmath,amssymb,amsthm,graphicx}
\usepackage{verbatim}
\newenvironment{reference}{\comment}{\endcomment}
%\newenvironment{reference}{}{}
\newenvironment{slogan}{\comment}{\endcomment}
\newenvironment{history}{\comment}{\endcomment}

% For commutative diagrams we use Xy-pic
\usepackage[all]{xy}

% We use 2cell for 2-commutative diagrams.
\xyoption{2cell}
% Theorem environments.
%
\theoremstyle{plain}
\newtheorem{theorem}[subsection]{Theorem}
\newtheorem{proposition}[subsection]{Proposition}
\newtheorem{lemma}[subsection]{Lemma}

\theoremstyle{definition}
\newtheorem{definition}[subsection]{Definition}
\newtheorem{example}[subsection]{Example}
\newtheorem{exercise}[subsection]{Exercise}
\newtheorem{situation}[subsection]{Situation}

\theoremstyle{remark}
\newtheorem{remark}[subsection]{Remark}
\newtheorem{remarks}[subsection]{Remarks}

\numberwithin{equation}{subsection}

% Macros
%
\def\lim{\mathop{\mathrm{lim}}\nolimits}
\def\colim{\mathop{\mathrm{colim}}\nolimits}
\def\Spec{\mathop{\mathrm{Spec}}}
\def\Hom{\mathop{\mathrm{Hom}}\nolimits}
\def\Ext{\mathop{\mathrm{Ext}}\nolimits}
\def\SheafHom{\mathop{\mathcal{H}\!\mathit{om}}\nolimits}
\def\SheafExt{\mathop{\mathcal{E}\!\mathit{xt}}\nolimits}
\def\Sch{\mathit{Sch}}
\def\Mor{\mathop{\mathrm{Mor}}\nolimits}
\def\Ob{\mathop{\mathrm{Ob}}\nolimits}
\def\Sh{\mathop{\mathit{Sh}}\nolimits}
\def\NL{\mathop{N\!L}\nolimits}
\def\CH{\mathop{\mathrm{CH}}\nolimits}
\def\proetale{{pro\text{-}\acute{e}tale}}
\def\etale{{\acute{e}tale}}
\def\QCoh{\mathit{QCoh}}
\def\Ker{\mathop{\mathrm{Ker}}}
\def\Im{\mathop{\mathrm{Im}}}
\def\Coker{\mathop{\mathrm{Coker}}}
\def\Coim{\mathop{\mathrm{Coim}}}

% Boxtimes
%
\DeclareMathSymbol{\boxtimes}{\mathbin}{AMSa}{"02}

%
% Macros for moduli stacks/spaces
%
\def\QCohstack{\mathcal{QC}\!\mathit{oh}}
\def\Cohstack{\mathcal{C}\!\mathit{oh}}
\def\Spacesstack{\mathcal{S}\!\mathit{paces}}
\def\Quotfunctor{\mathrm{Quot}}
\def\Hilbfunctor{\mathrm{Hilb}}
\def\Curvesstack{\mathcal{C}\!\mathit{urves}}
\def\Polarizedstack{\mathcal{P}\!\mathit{olarized}}
\def\Complexesstack{\mathcal{C}\!\mathit{omplexes}}
% \Pic is the operator that assigns to X its picard group, usage \Pic(X)
% \Picardstack_{X/B} denotes the Picard stack of X over B
% \Picardfunctor_{X/B} denotes the Picard functor of X over B
\def\Pic{\mathop{\mathrm{Pic}}\nolimits}
\def\Picardstack{\mathcal{P}\!\mathit{ic}}
\def\Picardfunctor{\mathrm{Pic}}
\def\Deformationcategory{\mathcal{D}\!\mathit{ef}}
\usepackage{cleveref}
\crefname{originaldefinition}{original definition}{original definitions}
\Crefname{originaldefinition}{Original definition}{Original definitions}
\makeatletter
\newlabel{categories-definition-directed}{{002V}{source}{}{sourcepointer-1}{}}
\newlabel{categories-definition-directed@cref}{{[originaldefinition][0][]002V}{[1][0][]source}{}{}{}}
\makeatother
\crefname{originallemma}{original lemma}{original lemmas}
\Crefname{originallemma}{Original lemma}{Original lemmas}
\makeatletter
\newlabel{lemma-almost-directed}{{00X4}{source}{}{sourcepointer-2}{}}
\newlabel{lemma-almost-directed@cref}{{[originallemma][0][]00X4}{[1][0][]source}{}{}{}}
\makeatother
\crefname{originalsection}{original section}{original sections}
\Crefname{originalsection}{Original section}{Original sections}
\makeatletter
\newlabel{categories-section-finite-limits}{{04AS}{source}{}{sourcepointer-3}{}}
\newlabel{categories-section-finite-limits@cref}{{[originalsection][0][]04AS}{[1][0][]source}{}{}{}}
\makeatother
\crefname{originallemma}{original lemma}{original lemmas}
\Crefname{originallemma}{Original lemma}{Original lemmas}
\makeatletter
\newlabel{categories-lemma-split-into-directed}{{002X}{source}{}{sourcepointer-4}{}}
\newlabel{categories-lemma-split-into-directed@cref}{{[originallemma][0][]002X}{[1][0][]source}{}{}{}}
\makeatother
\title{Conditions making the auxiliary index categories filtered}
\author{The Stacks Project Authors}
\date{stacks-project:sites.tex:local-snapshot-20261012; GNU FDL 1.2 or later}
\begin{document}
\maketitle
\noindent\textit{Editorial note.} All delivered mathematical text is the author's own, copied from the pinned Stacks Project chapter file sites-cohomology.tex; no statement or proof step was rewritten, reordered or removed.\par\medskip
\noindent\textit{Editorial note.} The proof is the author's complete argument; every step the author writes out is inside the counted span and no closing sentence is a placeholder.\par\medskip
\noindent\textit{Editorial note.} Every labelled reference occurring inside the counted spans is carried in assembly.external\_references as a non-load-bearing navigation pointer (same-chapter ones at their definition site, cross-chapter ones at the invocation inside the counted span), so no load-bearing mathematics is replaced by a stub. No bibliography citation occurs inside the counted spans.\par\medskip
\noindent\textit{Editorial note.} The checked chapter list for same-name collisions is the set of stacks chapters staged in this repository: chow, cohomology, divisors, etale-cohomology, flat, more-algebra, more-morphisms, morphisms, spaces-more-morphisms, varieties together with this card's chapter sites-cohomology.\par\medskip
\noindent\textit{Editorial note.} The preamble is the chapter's own retained minimal preamble (sources/stacks-sites-cohomology/retained-preamble.tex), a verbatim excerpt of the native preamble: the Xy-pic setup and the native theorem declarations with the mathematical macros.\par\medskip
\noindent\textit{Allowed standard prerequisites (editorial).}\begin{itemize}
\item Auxiliary categories are almost directed. Same chapter, Lemma lemma-almost-directed. Supplies the properties verified for the auxiliary categories.
\item Almost directed categories split into directed components. Categories, Lemma categories-lemma-split-into-directed. Reduces the assertion to connectedness.
\end{itemize}\medskip
\ifdefined\hypertarget\hypertarget{sourcepointer-1}{}\fi
\noindent\textit{Source reference.} Original definition 002V: sources/stacks-sites/pinned-original-source.tex, source lines 679--679. Navigation pointer invoked at proof line 679; the counted span contains the pointer only and the cited result is not reproduced inside it, so the pointer is registered as a non-load-bearing navigation stub in the established form for Stacks material. No bibliography citation is involved.\par\medskip
\ifdefined\hypertarget\hypertarget{sourcepointer-2}{}\fi
\noindent\textit{Source reference.} Original lemma 00X4: sources/stacks-sites/pinned-original-source.tex, source lines 684--684. Navigation pointer invoked at proof line 684; the counted span contains the pointer only and the cited result is not reproduced inside it, so the pointer is registered as a non-load-bearing navigation stub in the established form for Stacks material. No bibliography citation is involved.\par\medskip
\ifdefined\hypertarget\hypertarget{sourcepointer-3}{}\fi
\noindent\textit{Source reference.} Original section 04AS: sources/stacks-sites/pinned-original-source.tex, source lines 686--686. Navigation pointer invoked at proof line 686; the counted span contains the pointer only and the cited result is not reproduced inside it, so the pointer is registered as a non-load-bearing navigation stub in the established form for Stacks material. No bibliography citation is involved.\par\medskip
\ifdefined\hypertarget\hypertarget{sourcepointer-4}{}\fi
\noindent\textit{Source reference.} Original lemma 002X: sources/stacks-sites/pinned-original-source.tex, source lines 639--639. Navigation pointer invoked at proof line 639; the counted span contains the pointer only and the cited result is not reproduced inside it, so the pointer is registered as a non-load-bearing navigation stub in the established form for Stacks material. No bibliography citation is involved.\par\medskip
\begin{lemma}
\label{lemma-directed}
Let $u : \mathcal{C} \to \mathcal{D}$ be a functor between categories.
Assume
\begin{enumerate}
\item the category $\mathcal{C}$ has a final object $X$ and
$u(X)$ is a final object of $\mathcal{D}$ , and
\item the category $\mathcal{C}$ has fibre products and
$u$ commutes with them.
\end{enumerate}
Then the index categories $(\mathcal{I}^u_V)^{opp}$ are filtered (see
Categories, Definition \ref{categories-definition-directed}).
\end{lemma}

\begin{proof}
The assumptions imply that the assumptions of
Lemma \ref{lemma-almost-directed}
are satisfied (see the discussion in
Categories, Section \ref{categories-section-finite-limits}).
By
Categories, Lemma \ref{categories-lemma-split-into-directed}
we see that $\mathcal{I}_V$ is a (possibly empty) disjoint union of
directed categories.
Hence it suffices to show that $\mathcal{I}_V$ is connected.

\medskip\noindent
First, we show that $\mathcal{I}_V$ is nonempty.
Namely, let $X$ be the final object of $\mathcal{C}$,
which exists by assumption.
Let $V \to u(X)$ be the morphism coming from the fact
that $u(X)$ is final in $\mathcal{D}$ by assumption.
This gives an object of $\mathcal{I}_V$.

\medskip\noindent
Second, we show that $\mathcal{I}_V$ is connected.
Let $\phi_1 : V \to u(U_1)$ and $\phi_2 : V \to u(U_2)$ be
in $\Ob(\mathcal{I}_V)$. By assumption $U_1\times U_2$
exists and $u(U_1\times U_2) = u(U_1)\times u(U_2)$.
Consider the morphism $\phi : V \to u(U_1\times U_2)$
corresponding to $(\phi_1, \phi_2)$ by the universal property
of products. Clearly the object $\phi : V \to u(U_1\times U_2)$
maps to both $\phi_1 : V \to u(U_1)$ and $\phi_2 : V \to u(U_2)$.
\end{proof}

\end{document}
